\RequirePackage{pdfmanagement}
\SetKeys[document/metadata]{pdfversion=1.7, pdfstandard=A-2u, lang=en-GB}
\input{preamble/pdfA-compliance/glyphtounicode}
\documentclass[11pt,a4paper]{report}

\newcommand*{\PaperTitle}{Certified minimum numbers of matches for four heat-recovery benchmark instances}
\newcommand*{\PaperTitlePlain}{Certified minimum numbers of matches for four heat-recovery benchmark instances}
\newcommand*{\PaperShortTitle}{Certified minimum numbers of matches}
\newcommand*{\PaperSubtitle}{chemistry-open-problems \textperiodcentered{} 25 October 2026 \textperiodcentered{} solved (preliminary) for four instances, partial result for the others}
\newcommand*{\PaperCredits}{Solver: Claude Opus 5.5 (effort xhigh), sub-agent of the orchestrator, in two sessions (24 and 25 October 2026) \textperiodcentered{} Verification: pending}
\newcommand*{\PaperKeywords}{heat exchanger network synthesis; minimum number of matches; transshipment model; exact mixed-integer programming; VIPR certificate; verified computation}
\input{chemop-preamble}
\addbibresource{references.bib}

\newcommand*{\inst}[1]{\texttt{#1}}

\begin{document}
\makepapertitle

\begin{paperabstract}
Letsios, Kouyialis and Misener (2018) published 48 benchmark instances of the minimum number of matches problem,
the NP-hard mixed-integer linear program at the core of sequential heat exchanger network synthesis, and observed
that state-of-the-art solvers cannot solve moderately sized instances to global optimality. We determine the exact
optimum of four instances that were left without a proof of optimality: \inst{37sp-yfyv}${}=36$,
\inst{20sp1}${}=19$, \inst{23sp1}${}=22$ (best value known before: 23) and \inst{unbalanced10}${}=25$. Every upper
bound is an explicit network, checked in exact rational arithmetic by two separately written programs, the
second one against the original model of Letsios et al.\ with their own big-M values. The lower bounds of the first
three instances follow from a block lemma (every connected component of the match graph is heat balanced and
cascade feasible) and an exhaustive exact enumeration of blocks; the same argument raises the lower bound of the four
33-stream \inst{balanced15} instances from 31 to 32. The lower bound for \inst{unbalanced10} is a VIPR certificate
(0.54\,GB) of a pure LP-based branch-and-bound proof, written by exact-arithmetic SCIP~10.0.2 for a strengthened
transshipment model; every row of that model is proved valid here and re-derived separately from the data, and
the certificate is accepted both by VIPR's checker and by a second checker written from the format specification.
On the way we found that VIPR's reference checker accepts two kinds of invalid inference steps (rounding over
continuous variables, branching disjunctions with a fractional right-hand side); we give a patch and test
certificates.
\end{paperabstract}

\section*{The question}

The source is \textcite{LKM2018} (arXiv v2, Sec.~9.2, p.~30), whose data, code and solver logs are public:
\begin{openquote}
``Observe that state-of-the-art approaches cannot, in general, solve moderately-sized problems with 30-40 streams
to global optimality. For example, none of the test cases in the Grossmann (2017) or large scale test sets is
solved to global optimality within the specified time limit.''
\end{openquote}
The same section states that each run used ``absolute gap 0.99, relative gap 4\%''. The solver logs of the
repository show that \inst{unbalanced10} (value 25) and \inst{37sp-yfyv} (value 36) are among the instances
reported as solved, but every one of their six runs (CPLEX 12.6.3 and Gurobi 6.5.2, three formulations) stopped at
the 4\,\% tolerance with best bound 24, respectively 35. For \inst{20sp1} and \inst{23sp1} the time limit was
reached (best bounds 16 and 17). Twenty-five instances of the benchmark had no matching upper and lower bound.

\begin{theorem}\label{thm:main}
On the data files of the benchmark of \textcite{LKM2018}, the minimum number of matches is
\begin{center}
\inst{37sp-yfyv}: 36, \qquad \inst{20sp1}: 19, \qquad \inst{23sp1}: 22, \qquad \inst{unbalanced10}: 25.
\end{center}
Moreover the minimum number of matches is at least 32 for \inst{balanced15} and \inst{balanced15\_random0/1/2}
(previous bound 31), at least 32 for \inst{unbalanced15}, at least 23 for \inst{22sp1}, and at least 26 for
\inst{balanced12} and \inst{balanced12\_random0/1/2}. The optimum of the transshipment MILP of
\textcite{LKM2018}, with any of the three big-M files of their repository, has the same value as the minimum number
of matches (\zcref{prop:milp}).
\end{theorem}

\begin{table}[h]
\centering\footnotesize
\caption{Status of the 25 instances without matching bounds in 2018. LKM18: best upper / lower bound over the six
runs in the repository of \textcite{LKM2018} (floating-point solver output, lower bounds rounded up). Here: exact
network (UB) and \zcref{lem:block} or a VIPR certificate (LB). Families: \inst{b12} stands for \inst{balanced12} and
\inst{balanced12\_random0/1/2}, likewise \inst{b15}, \inst{u17} (\inst{unbalanced17}), \inst{u20}.}\label{tab:status}
\setlength{\tabcolsep}{4pt}
\begin{tabular}{@{}lcccl@{}}
\toprule
instance & $N$ & LKM18 UB / LB & here UB / LB & status \\
\midrule
\inst{37sp-yfyv} & 38 & 36 / 35 & 36 / 36 & optimal 36 (block lemma) \\
\inst{20sp1} & 21 & 19 / 16 & 19 / 19 & optimal 19 (block lemma) \\
\inst{23sp1} & 24 & 23 / 17 & 22 / 22 & optimal 22 (block lemma), new UB \\
\inst{unbalanced10} & 23 & 25 / 24 & 25 / 25 & optimal 25 (VIPR certificate) \\
\inst{unbalanced15} & 33 & 36 / 35 & 36 / 32 & open \\
\inst{22sp1} & 24 & 25 / 23 & 25 / 23 & open \\
\inst{b12} (4 instances) & 27 & 28--29 / 26--27 & 28--29 / 26 & open \\
\inst{b15} (4 instances) & 33 & 35 / 31 & 35 / 32 & open, new LB \\
\inst{u17} (4 instances) & 41 & 43 / 39--40 & 43 / -- & open \\
\inst{u20} (4 instances) & 47 & 50--51 / 45 & 50--51 / -- & open \\
\inst{large\_scale0/1/2} & 162 & 175--239 / 123--135 & -- & open, not attempted \\
\bottomrule
\end{tabular}
\end{table}


For \inst{balanced12\_random1} the certified bound 26 is below the floating-point bound 27 of the 2018 runs; no
certified lower bound in \zcref{tab:status} contradicts a floating-point one.

\section{The problem}

An instance (file \texttt{data/mip\_instances/<set>/<name>.dat} of the repository of \textcite{LKM2018}) has hot
streams $i\in H$ ($|H|=n$, hot utilities included with the loads fixed by the minimum utility cost step), cold
streams $j\in C$ ($|C|=m$, cold utilities included) and temperature intervals $t=0,\dots,k-1$, interval $0$ being
the hottest. $Q^H_{it}\ge0$ is the heat that hot stream $i$ supplies in interval $t$ and $Q^C_{jt}\ge0$ the heat that
cold stream $j$ needs in interval $t$. All numbers are read as exact decimal rationals. We write
$H_i[a..b]=\sum_{a\le t\le b}Q^H_{it}$, $H_i({\le}t)=H_i[0..t]$, $H_i=H_i({\le}k{-}1)$, and $D_j[a..b]$,
$D_j({\le}t)$, $D_j$ likewise; $N=n+m$; $R_t=\sum_{s<t}\bigl(\sum_iQ^H_{is}-\sum_jQ^C_{js}\bigr)$ is the heat cascade.
For all 54 data files we checked exactly that $R_k=0$, $R_t\ge0$ for all $t$, that the residuals printed in the file
equal $R_t$, and that every stream has positive total heat.

\begin{definition}\label{def:network}
A \emph{feasible network} is a family of numbers $q_{ijt}\ge0$ such that $\sum_iq_{ijt}=Q^C_{jt}$ for all $j,t$ and
$r_{it}:=H_i({\le}t)-\sum_{s\le t}\sum_jq_{ijs}\ge0$ for all $i,t$, with $r_{i,k-1}=0$. Its \emph{match set} is
$E(q)=\{(i,j):\sum_tq_{ijt}>0\}$, and the \emph{minimum number of matches} is $\min_q|E(q)|$.
\end{definition}

Heat of hot stream $i$ produced in interval $s$ can thus be used in interval $s$ or in any colder interval (the
minimum approach temperature is built into the intervals); $r_{it}$ is the heat of $i$ carried from interval $t$ to
interval $t+1$. The transshipment MILP of \textcite{PapouliasGrossmann1983}, as implemented by \textcite{LKM2018}
(\texttt{lib/exact\_methods/transshipment\_mip\_solver.py}), reads: minimise $\sum_{ij}y_{ij}$ subject to
$\sum_jq_{ijt}+\rho_{i,t+1}=Q^H_{it}+\rho_{it}$, $\sum_iq_{ijt}=Q^C_{jt}$, $\sum_tq_{ijt}\le U_{ij}y_{ij}$,
$\sum_i\rho_{i0}=\sum_i\rho_{ik}=0$, $y\in\{0,1\}$, $q,\rho\ge0$, with big-M values $U_{ij}$ read from the
files \texttt{<name>.bgm} of the directories \texttt{simple}, \texttt{gundersen} and \texttt{greedy} under
\texttt{data/bigM\_parameters/<set>/}.

\begin{proposition}\label{prop:milp}
For each instance of \zcref{thm:main} and each of the three big-M files, the optimum of this MILP equals the minimum
number of matches.
\end{proposition}
\begin{proof}
If $(y,q,\rho)$ is feasible for the MILP, then $\rho_{i,t+1}=r_{it}$, so $q$ is a feasible network and
$E(q)\subseteq\{(i,j):y_{ij}=1\}$; hence the MILP optimum is at least the minimum number of matches, whatever the
$U_{ij}$. Conversely, for each network of \zcref{sec:ub}, setting $y_{ij}=1$ exactly on $E(q)$ and
$\rho_{i,t+1}=r_{it}$ gives a feasible point: the script \texttt{check\_network\_lkm18.py} checks in exact
arithmetic every constraint of the MILP, including $\sum_tq_{ijt}\le U_{ij}y_{ij}$ for the three big-M files, and that
the objective equals the stated number of matches. With the lower bounds of \zcref{sec:block,sec:vipr} this gives
equality.
\end{proof}

\section{Upper bounds: explicit networks}\label{sec:ub}

The directory \texttt{certificates/solutions/} contains, for 22 instances, a network given by rational heats
$q_{ijt}$ (JSON). They were found with floating-point solvers and with \zcref{lem:structure} below, then rounded to
exact flows by an integer max-flow computation (all data scaled by the least common multiple of the decimal
denominators). Two programs check them:
\begin{itemize}
\item \texttt{check\_solutions.py} (with \texttt{hen\_exact.py}) checks $q\ge0$, every cold balance, $r_{it}\ge0$,
$r_{i,k-1}=0$, and that every declared match carries positive heat;
\item \texttt{check\_network\_lkm18.py}, written separately with its own parsers, checks the network as a point of
the MILP of \textcite{LKM2018} with $y$ set to the indicator of $E(q)$ (\zcref{prop:milp}).
\end{itemize}
Both print \texttt{ALL PASS} for the 22 networks, which have 36, 19, 22, 25 matches for \inst{37sp-yfyv},
\inst{20sp1}, \inst{23sp1}, \inst{unbalanced10}, and reproduce the best 2018 values for the other 18 instances of
\zcref{tab:status} (they improve none of them).

\section{The block lemma}\label{sec:block}

\begin{definition}
A nonempty set $K\subseteq H\cup C$ of streams is a \emph{block} if
$\sum_{i\in K\cap H}H_i=\sum_{j\in K\cap C}D_j$ and $\sum_{i\in K\cap H}H_i({\le}t)\ge\sum_{j\in K\cap C}D_j({\le}t)$
for every $t$. Let $c_{\max}$ be the largest number of parts of a partition of $H\cup C$ into blocks.
\end{definition}

\begin{lemma}\label{lem:block}
Every feasible network $q$ satisfies $|E(q)|\ge N-c_{\max}$.
\end{lemma}
\begin{proof}
Let $G$ be the bipartite graph on $H\cup C$ with edge set $E(q)$. Every stream has positive total heat, so it has an
incident edge; every connected component $K$ of $G$ therefore has at least $|K|-1\ge1$ edges. All heat of a hot
stream of $K$ goes to its neighbours, which lie in $K$, and all demand of a cold stream of $K$ comes from its
neighbours, so $\sum_{i\in K}H_i=\sum_{i,j\in K,t}q_{ijt}=\sum_{j\in K}D_j$, and for every $t$
\[\textstyle\sum_{j\in K}D_j({\le}t)=\sum_{j\in K}\sum_{s\le t}\sum_{i\in K}q_{ijs}
=\sum_{i\in K}\bigl(H_i({\le}t)-r_{it}\bigr)\le\sum_{i\in K}H_i({\le}t).\]
Hence $K$ is a block, the $c(G)$ components form a partition into blocks, $c(G)\le c_{\max}$, and
$|E(q)|\ge\sum_K(|K|-1)=N-c(G)\ge N-c_{\max}$.
\end{proof}

\begin{lemma}\label{lem:structure}
If $c_{\max}=2$ and $|E(q)|=N-2$, then $G$ is a forest with two components $A,B$, $\{A,B\}$ is a partition into two
blocks, and the restriction of $q$ to $A$ (to $B$) is a feasible network of the sub-instance $A$ ($B$) whose match
graph is a spanning tree.
\end{lemma}
\begin{proof}
In the proof of \zcref{lem:block}, $N-2=|E|\ge N-c(G)$ forces $c(G)=2$ and exactly $|K|-1$ edges per component.
\end{proof}
\zcref{lem:structure} only guides the search for networks; it is not used in any proof of a bound.

\paragraph{Computation of $c_{\max}$} (\texttt{block\_bound.py}). With $L$ the least common multiple of the decimal
denominators, stream $v$ gets the integer weight $w_v=LH_i$ (hot) or $-LD_j$ (cold) and an integer cumulative
profile. All zero-weight subsets are enumerated exactly by meet in the middle (all subset sums of the two halves,
exact hash join); a zero-weight subset is a block if and only if the sum of the two half-profiles is nonnegative in
every interval (all partial sums are below $2^{62}$, asserted). A 2-block partition exists iff some block $A$
containing stream $0$, $A\ne H\cup C$, has a block as complement. A partition into three blocks exists iff there are
blocks $A\ni0$ and $B$ such that $X=(H\cup C)\setminus A$ is a block, $B\subsetneq X$ contains the lowest-index stream
of $X$, and $X\setminus B$ is a block (unions and differences of nested blocks are blocks, since both defining
conditions are additive); the search over all such pairs is exhaustive. A partition into four or more blocks yields
one into three by merging, so ``no 3-block partition'' means $c_{\max}\le2$.

\begin{table}[h]
\centering\footnotesize
\caption{Exhaustive exact enumeration (\texttt{run\_block\_bounds.sh}, 41\,s in total). ``2-part.'': number of
partitions into two blocks; no instance has a partition into three blocks.}\label{tab:block}
\setlength{\tabcolsep}{4pt}
\begin{tabular}{@{}lrrrrc@{}}
\toprule
instance & $N$ & zero-sum subsets & blocks & 2-part. & $N-c_{\max}$ \\
\midrule
\inst{37sp-yfyv} & 38 & 4\,833 & 2\,431 & 14 & 36 \\
\inst{20sp1} & 21 & 15 & 12 & 4 & 19 \\
\inst{23sp1} & 24 & 17 & 12 & 3 & 22 \\
\inst{22sp1} & 24 & 1 & 1 & 0 & 23 \\
\inst{balanced12}, \inst{\_random0/1/2} & 27 & $\approx6\cdot10^4$ & $\approx1.8\cdot10^4$ & 0 & 26 \\
\inst{balanced15} & 33 & 1\,817\,217 & 518\,054 & 0 & 32 \\
\inst{balanced15\_random0/1/2} & 33 & $1.75$--$3.42\cdot10^6$ & $0.63$--$1.00\cdot10^6$ & 0 & 32 \\
\inst{unbalanced15} & 33 & 470\,203 & 94\,888 & 0 & 32 \\
\bottomrule
\end{tabular}
\end{table}


With the networks of \zcref{sec:ub} this proves the values 36, 19 and 22 of \zcref{thm:main}. For \inst{unbalanced10}
($N=23$) the block lemma gives at most 22, so a different argument is needed.

\section{unbalanced10: a strengthened model and a VIPR certificate}\label{sec:vipr}

\subsection{The strengthened transshipment model}

Let $\mathrm{rc}_i(t)=\min\bigl(R_{t+1},H_i({\le}t)\bigr)$ for $0\le t\le k-2$ and
$\mathrm{rc}_i(-1)=\mathrm{rc}_i(k-1)=0$.

\begin{lemma}\label{lem:rc}
In every feasible network, $0\le r_{it}\le\mathrm{rc}_i(t)$.
\end{lemma}
\begin{proof}
Summing the definition of $r_{it}$ over $i$ and using the cold balances gives
$\sum_ir_{it}=\sum_iH_i({\le}t)-\sum_jD_j({\le}t)=R_{t+1}$; all $r_{i't}$ are nonnegative, so $r_{it}\le R_{t+1}$,
and $r_{it}\le H_i({\le}t)$ by definition.
\end{proof}

Given a feasible network $q$, put $y_{ij}=1$ if $(i,j)\in E(q)$ and $y_{ij}=0$ otherwise.

\begin{proposition}\label{prop:rows}
$(q,r,y)$ satisfies each of the following constraints. Consequently, if the MILP formed by them together with
$\sum y\le K$ has no solution, every feasible network has at least $K+1$ matches.
\begin{enumerate}
\item Balances: $\sum_jq_{ijt}+r_{it}-r_{i,t-1}=Q^H_{it}$ and $\sum_iq_{ijt}=Q^C_{jt}$ (with $r_{i,-1}=r_{i,k-1}=0$).
\item Bounds: $0\le r_{it}\le\mathrm{rc}_i(t)$ and $0\le q_{ijt}\le u_{ijt}:=\min\bigl(Q^C_{jt},Q^H_{it}+\mathrm{rc}_i(t-1)\bigr)$.
\item Disaggregated links: $q_{ijt}\le u_{ijt}y_{ij}$.
\item Aggregated links: $\sum_tq_{ijt}\le U_{ij}y_{ij}$, where $U_{ij}$ is the minimum over all labellings
$\lambda:\{0,\dots,k-1\}\to\{\mathrm{in},\mathrm{out}\}$ of
$\mathrm{cap}(\lambda)=\sum_{\lambda(u)=\mathrm{in}}Q^C_{ju}+\sum_{\lambda(u)=\mathrm{out}}Q^H_{iu}
+\sum_{\lambda(u)=\mathrm{in},\lambda(u+1)=\mathrm{out}}\mathrm{rc}_i(u)$.
\item Cold windows ($Q^C_{jT}>0$, $s\le T$):
\[\textstyle\sum_i\min\bigl(D_j[s..T],\,\mathrm{rc}_i(s-1)+H_i[s..T]\bigr)y_{ij}\ge D_j[s..T].\]
\item Hot windows ($Q^H_{is}>0$, $s\le T$, right-hand side positive):
\[\textstyle\sum_j\min\bigl(H_i[s..T],D_j[s..T]\bigr)y_{ij}\ge H_i[s..T]-\mathrm{rc}_i(T).\]
\item Partner counts: for each window row $\sum_ea_ey_e\ge b$, $\sum_ey_e\ge p$, where $p$ is the least number of
coefficients, taken largest first, whose sum reaches $b$.
\end{enumerate}
\end{proposition}
\begin{proof}
(1) are the definitions. (2) $r$: \zcref{lem:rc}. $q$: the heat that $i$ delivers in interval $t$ is at most
$r_{i,t-1}+Q^H_{it}$ and the heat that $j$ receives there is $Q^C_{jt}$. (3) is (2) when $y_{ij}=1$, and
$q_{ijt}=0$ when $y_{ij}=0$. (4) Fix a labelling $\lambda$. Intervals labelled in contribute $q_{iju}\le Q^C_{ju}$. On
a maximal run $[a..b]$ of out intervals, the heat $i$ delivers to anybody is
$r_{i,a-1}+H_i[a..b]-r_{ib}\le\mathrm{rc}_i(a-1)+H_i[a..b]$, and $\mathrm{rc}_i(a-1)$ is paid by the in-to-out
transition at $a-1$ (or vanishes if $a=0$). So $\sum_tq_{ijt}\le\mathrm{cap}(\lambda)$, and $=0$ if $y_{ij}=0$.
(5) $D_j[s..T]=\sum_i\sum_{s\le t\le T}q_{ijt}$, and the $i$-th inner sum is at most $D_j[s..T]$, at most the heat
$i$ delivers in the window, $r_{i,s-1}+H_i[s..T]-r_{iT}\le\mathrm{rc}_i(s-1)+H_i[s..T]$, and zero if $y_{ij}=0$.
(6) Decompose the flow along the chain of hot stream $i$ into amounts $g(a\to t)\ge0$, $a\le t$, of heat produced in
$a$ and delivered in $t$ ($\sum_tg(a\to t)=Q^H_{ia}$, $\sum_ag(a\to t)=\sum_jq_{ijt}$, e.g.\ first in, first out),
and split each delivery among the cold streams in proportion to $q_{ijt}$. Let $O_j$ be the heat produced in
$[s..T]$ and delivered to $j$ inside $[s..T]$: $O_j\le\min(H_i[s..T],D_j[s..T])\,y_{ij}$. Heat produced in $[s..T]$
and delivered after $T$ is part of $r_{iT}$, so $\sum_jO_j\ge H_i[s..T]-r_{iT}\ge H_i[s..T]-\mathrm{rc}_i(T)$.
(7) With binary $y$, a knapsack row $\sum a_ey_e\ge b$, $a_e\ge0$, needs at least $p$ variables equal to one.
Finally, pairs with $U_{ij}=0$ and variables with $u_{ijt}=0$ or $\mathrm{rc}_i(t)=0$ are omitted from the model;
they are zero in every feasible network by (2) and (4).
\end{proof}

The minimum in (4) is computed by a two-state dynamic programme (\texttt{hen\_cuts.py}). It is the capacity of a
minimum cut in the two-stream network source${}\to{}$hot$_u$ ($Q^H_{iu}$), hot$_u\to{}$hot$_{u+1}$
($\mathrm{rc}_i(u)$), hot$_u\to{}$cold$_u$ ($\infty$), cold$_u\to{}$sink ($Q^C_{ju}$), but validity only uses the
direct argument above.

\subsection{The certified computation}

The model of \zcref{prop:rows} for \inst{unbalanced10}, with objective $\sum y$ and the extra row
\texttt{maxm}: $\sum y\le24$, is the CPLEX-LP file \texttt{certificates/vipr/unbalanced10\_le24.lp} (122 binary,
816 + 147 continuous variables, 2453 rows of which 1247 are window and partner-count rows; exact decimal
coefficients; SHA-256 \texttt{e3dcb597\ldots}). Its certification has five steps, each printing PASS or FAIL
(\texttt{verify\_vipr.sh}):
\begin{enumerate}
\item the LP file is regenerated from the data by \texttt{hen\_lpfile.py} and must be byte-identical;
\item \texttt{check\_lp\_model.py} re-derives every row and bound of the LP file from the data with separately
written code (exact cascade recomputed from the data, $U_{ij}$ checked as the capacity of an explicitly
reconstructed labelling, window coefficients, partner counts, the row \texttt{maxm}), checks that every omitted
variable is forced to zero, and that the file has no other bounds or binaries;
\item \texttt{check\_vipr\_model.py} checks that every constraint of the certificate's problem section equals, up to a
positive factor, a row or a variable bound of the LP file, that every integer variable of the certificate is a
binary of the LP file, and that the relation to prove is \texttt{infeas}; the certificate therefore proves the
infeasibility of a relaxation of the LP-file MILP;
\item \texttt{viprchk} (VIPR, commit 30f2951, with the two patches of \zcref{sec:checkers}) verifies the certificate in
exact GMP rational arithmetic, on the raw file and on the output of \texttt{viprcomp};
\item \texttt{vipr\_check\_independent.py}, a second checker written from the format specification
\autocite{CheungGleixnerSteffy2017,VIPRtools} without using \texttt{viprchk}'s code, verifies the same file.
\end{enumerate}

The certificate was produced by SCIP 10.0.2 with SoPlex 8.0.2 \autocite{SCIP1002,EiflerGleixner2023}, built from
source with exact rational solving, run with \texttt{exact/enable = TRUE}, presolving, cutting planes and conflict
analysis switched off (\texttt{certificates/vipr/scip\_exact\_safe.set}), single thread. SCIP reports
infeasibility after 10\,233 nodes (5\,117 infeasible leaves, maximal depth 32; 287\,s on an idle machine, 553\,s
with other jobs running). The certificate (VIPR 1.0 syntax) has 540\,587\,890 bytes and SHA-256
\begin{center}\small\texttt{17967e67e607271c2fc60fc77a185da2badcbb84459d7deaeaddd5a4601ec352};\end{center} it has 1\,085 variables (122 integer), 4\,623 problem constraints
(2\,453 rows and 2\,170 variable bounds) and 191\,697 derivation steps: 10\,232 assumptions, 170\,552 linear
combinations (14.8 million multipliers in all), 5\,797 roundings and 5\,116 unsplitting (branching) steps; there is no
``weak'' or ``incomplete'' step. It is not stored in the repository: \texttt{run\_vipr\_certificate.sh} regenerates it
(the same 10\,233 nodes were obtained on 24 and 25 October with two different builds); its size and hash, the solver
log and the verification log are in \texttt{certificates/vipr/}.

\begin{center}\small
\begin{tabular}{llr}
\toprule
step & check & result \\
\midrule
1 & LP file regenerated from the data, byte-identical & PASS \\
2 & every row and bound re-derived from the data & PASS \\
3 & certificate problem $\subseteq$ LP model, RTP \texttt{infeas} & PASS \\
4a & patched \texttt{viprchk}, raw certificate & PASS (68\,s) \\
4b & \texttt{viprcomp}, then patched \texttt{viprchk} & PASS (64\,s) \\
5 & \texttt{vipr\_check\_independent.py} & PASS (199\,s) \\
\bottomrule
\end{tabular}
\end{center}

The whole of \texttt{verify\_vipr.sh} took 7\,min\,40\,s. The certificate is also accepted by the unpatched
\texttt{viprchk} and by the build with the canonicalisation patch only, so it uses neither of the patterns of
\zcref{sec:checkers}. As negative tests, \texttt{vipr\_mutation\_tests.sh} changes the certificate in four places
(right-hand side of $\sum y\le24$ replaced by 25, so that the problem becomes feasible; one multiplier changed; the
two assumptions of an unsplitting step swapped; a rounding step claiming $y\ge2$); both checkers reject each of the
four corrupted files. Together with the 25-match network of \zcref{sec:ub} this proves \inst{unbalanced10}${}=25$.

\paragraph{Trusted base.} The proofs of \zcref{lem:rc,prop:rows}; \texttt{check\_lp\_model.py} and
\texttt{check\_vipr\_model.py}; and either \texttt{viprchk} with its patches and GMP, or the second checker
(Python with gmpy2); the C++ and Python tool chains. SCIP, SoPlex, \texttt{viprcomp} and
the LP generator are not trusted.

\section{Checking the checkers}\label{sec:checkers}

\paragraph{Two gaps in VIPR's reference checker.} Reading \texttt{viprchk.cpp} (commit 30f2951; the current master,
commit 99eec19 of 24 October 2026, has the same code) we found two inference rules that are checked too weakly:
\begin{enumerate}
\item a rounding step (\texttt{rnd}) requires integer coefficients on the integer variables but does not reject
nonzero coefficients on continuous variables, although Chv\'atal--Gomory rounding is only valid when the whole
support is integer;
\item an unsplitting step (\texttt{uns}) checks that the two branching constraints are $ax\le\beta$ and
$ax\ge\beta+1$ with $a$ integral on integer variables, but not that $\beta$ is an integer; for fractional $\beta$
the disjunction misses integer points.
\end{enumerate}
The certificates \texttt{certificates/vipr/viprchk\_tests/bad\_rnd\_continuous.vipr} (``proves'' that
$x\ge\frac12$, $x\le\frac12$, $x$ continuous, is infeasible) and \texttt{bad\_uns\_fractional.vipr} (``proves'' that
$x\ge1$, $x\le1$, $x$ integer, is infeasible by branching on $x\le\frac12\mid x\ge\frac32$) are accepted by the
unpatched \texttt{viprchk}. The patch \texttt{certificates/viprchk\_harden.patch} adds both checks, and also bounds
checks on constraint indices (the original performs them only in debug builds); the patched checker and
\texttt{vipr\_check\_independent.py} reject both test certificates, and still accept the example certificates of the VIPR distribution that the original accepts. The
hardening patch only adds rejections.

\paragraph{Canonicalisation patch.} GMP's C++ stream input does not canonicalise rationals (``5/5'' is stored as
5/5, and GMP's equality test assumes canonical operands), while SCIP prints non-canonical fractions. The patch
\texttt{certificates/viprchk\_canonicalize.patch} canonicalises every rational after it is read (four places);
canonicalisation does not change any value.

\paragraph{Earlier SCIP certificates that failed.} With cutting planes at all nodes, or at the root only, SCIP
10.0.2 wrote Gomory mixed-integer cut derivations that \texttt{viprcomp} and \texttt{viprchk} could not complete
and verify (``Failed to derive constraint \ldots{}\_gom\ldots''). With conflict analysis on, a certificate of 2.19\,GB
contained 9\,828 ``weak'' conflict rows such as \texttt{AggrRow\_5376}, which claims that 116 of the 122 variables
sum to at most 18, derived from $\sum y\le24$ with the instruction to complete it with \emph{global} bounds; the row is
valid only under the local fixings of the six missing variables to one, so the derivation as written is wrong and the
verification failed at that row. We therefore use no cutting planes and no conflict analysis; with these settings
the certificate contains no weak derivation at all.

\paragraph{Floating-point pitfalls.} The 36-match network returned by floating-point SCIP for \inst{37sp-yfyv} is
infeasible in exact arithmetic (deficit exactly 0.1 heat units), and every one of the six match lists of the 2018
repository for this instance needs 37 pairs exactly after exact pruning (\texttt{side\_findings.py}); an exact
36-match network exists nonetheless (two spanning trees on a 2-block partition, \zcref{lem:structure}). The solution
writer of the 2018 code prints $y=1$ for every $y>10^{-10}$, so its match lists may contain pairs that carry
negligible heat; we always recompute flows exactly.

\section{Novelty, other instances, limitations}

\paragraph{Novelty check (25 October 2026).} The 13 works citing \textcite{LKM2018} in OpenAlex (and the 11 in
Semantic Scholar) were screened by title and abstract; the doctoral thesis of \textcite{Kouyialis2019} contains the
2018 study with the same values, and the survey \textcite{Letsios2020} reports no new values. The repositories of the
2018 study have no commit after February 2019; a GitHub search on the instance names and key phrases found only
those repositories and a student replication without bound tables; the AI-driven campaigns listed in the protocol of
this repository (GitHub users \texttt{ipitchford} and \texttt{mmislan}) have no heat-exchanger-network entry. To our
knowledge, none of the four optimal values had been proved before. \textcite{FurmanSahinidis2004} (read in full)
report CPLEX~7.0 values 19, 25, 23 and 36 for \inst{20sp1}, \inst{22sp1}, \inst{23sp1} and \inst{37sp-yfyv}, all
marked as not converged after 7\,h, and give no lower bounds; the value 22 for \inst{23sp1} is below the best value
reported there and in 2018 (23). \textcite{CerdaCafaro2022} (read in full) treat a different problem (detailed network
structures with stream splitting, 45 literature examples) and report no value for the data files used here.
\textcite{ChenGrossmannMiller2015}, the origin of the balanced and unbalanced instances, was not read in full.

\paragraph{Other instances.} We also tried the strengthened model on the open instances with the smallest gaps.
Floating-point SCIP 10.0.2 with its default settings (cutting planes on), 1800\,s per run, did not decide
$\sum y\le35$ for \inst{unbalanced15} (final dual bound 33.9), $\sum y\le24$ for \inst{22sp1} (21.5), or
$\sum y\le27$ for \inst{balanced12}, \inst{\_random0} and \inst{\_random2} ($\le28$ for \inst{\_random1}; dual
bounds 24.7--25.0). Without cutting planes, as required here for a
verifiable certificate, the root LP bounds of the strengthened model are 28.9 for \inst{unbalanced15} (36 needed),
18.2 for \inst{22sp1} (25 needed) and 22.2 for \inst{balanced12} (28 needed), against 21.1 (25 needed) for
\inst{unbalanced10}, which already took 10\,233 nodes; we did not run the certified pipeline on them. These
statements are exploration, not claims.

\paragraph{Limitations.} The data files are taken as given (decimal numbers read as exact); the underlying stream
temperatures and the minimum utility cost step were not re-derived. The lower bound for \inst{unbalanced10} rests on
software (\texttt{viprchk} or our second checker) and on the model-match scripts; it has not been re-checked by a
verifier yet.

\section{Reproduction}

All scripts are in \texttt{scripts/} and print PASS/FAIL; times are wall-clock on WSL2 (Fedora 44), one core of a
desktop CPU.
\begin{center}\footnotesize
\begin{tabular}{@{}lll@{}}
\toprule
script & checks & time \\
\midrule
\texttt{check\_solutions.py} & 22 networks, exact & 2\,s \\
\texttt{check\_network\_lkm18.py} & 22 networks as points of the 2018 MILP & 2\,s \\
\texttt{run\_block\_bounds.sh} & \zcref{lem:block} bounds, exact enumeration & 41\,s \\
\texttt{side\_findings.py} & floating-point findings for \inst{37sp-yfyv} & 17\,s \\
\texttt{run\_vipr\_certificate.sh} & produces the certificate (exact SCIP) & 5--10\,min \\
\texttt{verify\_vipr.sh} & steps 1--5 of \zcref{sec:vipr} & 7.7\,min \\
\texttt{vipr\_mutation\_tests.sh} & four corrupted certificates rejected & 6\,min \\
\bottomrule
\end{tabular}
\end{center}

\paperreferences
\end{document}
